\documentclass[border=10pt]{standalone}
\usepackage{tikz,ctex}
\usetikzlibrary{positioning}

\begin{document}

\begin{tikzpicture}[
    % 全局设置
    node distance = 1.6cm and 1.4cm,
    every node/.style = {
        draw,
        rounded corners,
        thick,
        align = center,
        font = \sffamily\small,
        minimum width = 2.0cm,
        minimum height = 0.9cm
    },
    % 样式定义
    centerStyle/.style = {
        fill = blue!45,
        text = white,
        font = \sffamily\bfseries\large,
        line width = 2pt,
        minimum size = 2.5cm,
        circle
    },
    branchTitle/.style = {
        font = \sffamily\bfseries,
        line width = 1.5pt
    },
    subNode/.style = {
        fill = white,
        font = \sffamily\small,
        minimum width = 2.0cm,
        minimum height = 0.8cm
    },
    leafNode/.style = {
        fill = white,
        font = \sffamily\footnotesize,
        minimum width = 1.8cm,
        minimum height = 0.7cm
    },
    arrowStyle/.style = {
        ->,
        thick,
        >=stealth,
        shorten >=2pt
    }
]

    % ================= 1. 中心节点 =================
    \node[centerStyle] (center) {7.4.归一化 \\ Normalization};

    % ================= 2. 上方分支：输入数据归一化 (红) =================
    \node[branchTitle, fill=red!15, above=of center] (data) {输入数据归一化 \\ Data Normalization};

    \node[subNode, above=of data, xshift=-6cm] (range) {量纲差异 \\ Different Ranges};
    \node[subNode, above=of data, xshift=-2cm] (mean) {均值 $\mu_i$ \\ Mean};
    \node[subNode, above=of data, xshift=2cm] (var) {方差 $\sigma_i^2$ \\ Variance};
    \node[subNode, above=of data, xshift=6cm] (rescale) {零均值单位方差 \\ Zero Mean Unit Var};

    % ================= 3. 右侧分支：批归一化 (蓝) =================
    \node[branchTitle, fill=orange!15, right=of center] (batch) {批归一化 \\ Batch Normalization};

    \node[subNode, right=of batch, yshift=4cm] (mini) {小批量统计 \\ Mini-batch Stats};
    \node[subNode, right=of batch, yshift=2cm] (preact) {预激活归一化 \\ Pre-activation Norm};
    \node[subNode, right=of batch, yshift=0cm] (gamma) {缩放参数 $\gamma_i$ \\ Scale};
    \node[subNode, right=of batch, yshift=-2cm] (beta) {平移参数 $\beta_i$ \\ Shift};
    \node[subNode, right=of batch, yshift=-4cm] (moving) {移动平均 \\ Moving Averages};

    % ================= 4. 下方分支：层归一化 (绿) =================
    \node[branchTitle, fill=green!15, below=of center] (layer) {层归一化 \\ Layer Normalization};

    \node[subNode, below=of layer, xshift=-6cm] (hidden) {跨隐藏单元 \\ Across Hidden Units};
    \node[subNode, below=of layer, xshift=-2cm] (perdata) {每数据点归一化 \\ Per Data Point};
    \node[subNode, below=of layer, xshift=2cm] (rnn) {RNN 与 Transformer \\ RNN \& Transformer};
    \node[subNode, below=of layer, xshift=6cm] (noavg) {无需移动平均 \\ No Moving Average};

    % ================= 5. 左侧分支：梯度问题与理论 (紫) =================
    \node[branchTitle, fill=purple!15, left=of center] (theory) {梯度问题与理论 \\ Gradient Issues \& Theory};

    \node[subNode, left=of theory, yshift=4cm] (vanish) {梯度消失 \\ Vanishing Gradients};
    \node[subNode, left=of theory, yshift=2cm] (explode) {梯度爆炸 \\ Exploding Gradients};
    \node[subNode, left=of theory, yshift=0cm] (smooth) {误差曲面平滑 \\ Smooth Error Surface};
    \node[subNode, left=of theory, yshift=-2cm] (covariate) {内部协变量偏移 \\ Internal Covariate Shift};
    \node[subNode, left=of theory, yshift=-4cm] (santurkar) {Santurkar 研究 \\ Santurkar et al.};

    % ================= 连线逻辑 =================
    % 中心 -> 四大分支标题
    \draw[arrowStyle, red] (center) -- (data);
    \draw[arrowStyle, blue] (center) -- (batch);
    \draw[arrowStyle, green!60!black] (center) -- (layer);
    \draw[arrowStyle, purple] (center) -- (theory);

    % 分支标题 -> 子节点 (上方)
    \draw[arrowStyle, red!60] (data) -- (range.south);
    \draw[arrowStyle, red!60] (data) -- (mean.south);
    \draw[arrowStyle, red!60] (data) -- (var.south);
    \draw[arrowStyle, red!60] (data) -- (rescale.south);

    % 分支标题 -> 子节点 (右侧)
    \draw[arrowStyle, blue!60] (batch) -- (mini.west);
    \draw[arrowStyle, blue!60] (batch) -- (preact.west);
    \draw[arrowStyle, blue!60] (batch) -- (gamma.west);
    \draw[arrowStyle, blue!60] (batch) -- (beta.west);
    \draw[arrowStyle, blue!60] (batch) -- (moving.west);

    % 分支标题 -> 子节点 (下方)
    \draw[arrowStyle, green!60!black] (layer) -- (hidden.north);
    \draw[arrowStyle, green!60!black] (layer) -- (perdata.north);
    \draw[arrowStyle, green!60!black] (layer) -- (rnn.north);
    \draw[arrowStyle, green!60!black] (layer) -- (noavg.north);

    % 分支标题 -> 子节点 (左侧)
    \draw[arrowStyle, purple!60] (theory) -- (vanish.east);
    \draw[arrowStyle, purple!60] (theory) -- (explode.east);
    \draw[arrowStyle, purple!60] (theory) -- (smooth.east);
    \draw[arrowStyle, purple!60] (theory) -- (covariate.east);
    \draw[arrowStyle, purple!60] (theory) -- (santurkar.east);

\end{tikzpicture}

\end{document}